\SourcePage{627}{630}
\section{The semiclassical case}

We shall trace how the classical limit emerges from the quantum-mechanical theory of scattering.

Excluding the zero scattering angle $\theta$ from consideration, we can write the scattering amplitude given by the exact theory in the form (124.4):
\begin{equation}
 f(\theta)=\frac1{2ik}\sum_{l=0}^{\infty}(2l+1)P_l(\cos\theta)e^{2i\delta_l}.
 \tag{127.1}
\end{equation}
We know that semiclassical wave functions have large phases. It is therefore natural to anticipate that the limiting transition in scattering theory corresponds to large phases $\delta_l$. The main contribution to the sum (127.1) comes from terms with large $l$. We may consequently replace $P_l(\cos\theta)$ by the asymptotic expression (49.7), written as
\[
 P_l(\cos\theta)\approx-\frac{i}{\sqrt{2\pi l\sin\theta}}
 \left\{\exp\left[i\left(l+\frac12\right)\theta+\frac{i\pi}4\right]
 -\exp\left[-i\left(l+\frac12\right)\theta-\frac{i\pi}4\right]\right\}.
\]
Substitution in (127.1) gives
\begingroup\small
\begin{equation}
 f(\theta)=\frac1k\sum_l\sqrt{\frac{l}{2\pi\sin\theta}}
 \left(\exp\left\{i\left[2\delta_l-\left(l+\frac12\right)\theta-\frac\pi4\right]\right\}
 -\exp\left\{i\left[2\delta_l+\left(l+\frac12\right)\theta+\frac\pi4\right]\right\}\right).
\tag{127.2}
\end{equation}
\endgroup

Regarded as functions of $l$, the exponential factors oscillate rapidly, since their phases are large. Most of the terms in (127.2) therefore cancel one another. The sum is determined chiefly by the range of $l$ close to the value at which the exponent of one of the exponentials is stationary, that is, close to a root of
\begin{equation}
 2\frac{d\delta_l}{dl}\pm\theta=0.
 \tag{127.3}
\end{equation}
This range contains many terms for which the exponential factors remain nearly constant, because their exponents vary slowly near the stationary point; these terms therefore do not cancel.

\SourcePage{628}{631}
In the semiclassical case, the phases $\delta_l$ can be written (see \S~124) as the limit, as $r\to\infty$, of the difference between the phase
\[
 \frac\pi4+\frac1\hbar\int_{r_0}^{r}
 \sqrt{2m[E-U(r)]-\frac{\hbar^2(l+1/2)^2}{r^2}}\,dr
\]
of the semiclassical wave function in the field $U(r)$ and the free-motion wave-function phase $kr-\pi l/2$ (see \S~33). Thus,
\begin{equation}
 \delta_l=\int_{r_0}^{\infty}\left[\sqrt{\frac{2m}{\hbar^2}(E-U)-\frac{(l+1/2)^2}{r^2}}-k\right]dr
 +\frac\pi2\left(l+\frac12\right)-kr_0.
 \tag{127.4}
\end{equation}
This expression is to be substituted into equation (127.3). When differentiating the integral, one must bear in mind that its limit $r_0$ also depends on $l$. The resulting term $k\,dr_0/dl$, however, is cancelled by the derivative of the term $-kr_0$ in $\delta_l$. The quantity $\hbar(l+1/2)$ is the particle's angular momentum. In classical mechanics it may be written as $m\rho v$, where $\rho$ is the \emph{impact parameter} and $v$ is the particle's velocity at infinity.


After making this substitution, equation (127.3) takes the final form
\begin{equation}
 2\int_{r_0}^{\infty}\frac{\rho\,dr}{r^2\sqrt{1-U/E-\rho^2/r^2}}=\pi\pm\theta.
 \tag{127.5}
\end{equation}
In a repulsive field this equation has a root for $\rho$ only with the minus sign before $\theta$ on the right-hand side, and in an attractive field only with the plus sign.

Equation (127.5) coincides exactly with the classical equation relating the scattering angle to the impact parameter (see I, \S~18). It is also straightforward to verify that the cross-section takes its classical form.

To do so, expand the exponent in (127.2) in powers of $l'=l-l_0(\theta)$, where $l_0(\theta)$ is determined by (127.3)--(127.5). For definiteness, consider the first term in (127.2), and accordingly take the lower sign in (127.3), the repulsive case. From (127.3), observe that
\[
 \frac{d\delta_l}{dl}=\frac12\frac{d\theta}{dl_0}.
\]

\SourcePage{629}{632}
We then have
\[
 i\left[2\delta_l-\left(l+\frac12\right)\theta-\frac\pi4\right]
 \approx i\left[2\delta_{l_0}-\left(l_0+\frac12\right)\theta-\frac\pi4\right]
 +i\frac12\frac{d\theta}{dl_0}l'^2.
\]
We now replace the summation over $l$ in (127.2) by integration over $dl'$ near $l'=0$. Treating $l'$ as a complex variable, direct the integration path near that point along the direction of steepest descent of the exponential, at an angle $\pi/4$ or $-\pi/4$ to the real axis according to the sign of $d\theta/dl_0$. In other words, put $l'=\xi\exp(\pm i\pi/4)$ and integrate over real $\xi$. Since the integral converges rapidly, its limits may be extended from $-\infty$ to $\infty$:
\[
 \int_{-\infty}^{\infty}\exp\left(-\frac12\left|\frac{d\theta}{dl_0}\right|\xi^2\right)d\xi
 =\sqrt{2\pi\left|\frac{dl_0}{d\theta}\right|}.
\]
The result is
\begin{equation}
 f(\theta)=\frac1k\left(\frac{l_0}{\sin\theta}\left|\frac{dl_0}{d\theta}\right|\right)^{1/2}
 \exp\left\{i\left[2\delta_{l_0}-\left(l_0+\frac12\right)\theta-\frac\pi4\right]\right\}.
 \tag{127.6}
\end{equation}
It follows that
\begin{equation}
 d\sigma=|f|^2\,2\pi\sin\theta\,d\theta
 =2\pi\frac{l_0}{k^2}\left|\frac{dl_0}{d\theta}\right|d\theta,
 \tag{127.7}
\end{equation}
and on introducing the impact parameter by $\rho=l_0/k$, we recover the classical formula $d\sigma=2\pi\rho\,d\rho$.

Thus, for scattering through a prescribed angle $\theta$ to be classical, both the value of $l$ satisfying (127.3) and the phase $\delta_l$ at that value must be large\footnote{The relation between $\theta$ and $\rho$ given by (127.5) need not be single-valued; more than one value of $\rho$ may then correspond to the same $\theta$. In that event, $f(\theta)$ is a sum of expressions (127.6) with the corresponding values of $l_0$. At extrema of $\theta(\rho)$, the derivative $d\rho/d\theta$, and with it the classical differential cross-section $d\sigma/do$, becomes infinite; the classical approximation is, of course, inadequate near such an angle (see Problem 2).}. This condition has a simple meaning. To speak of classical scattering through an angle $\theta$ for a particle passing with impact parameter $\rho$, the quantum-mechanical uncertainties in both quantities must be relatively small: $\Delta\rho\ll\rho$ and $\Delta\theta\ll\theta$. The uncertainty of the scattering angle is of order $\Delta\theta\sim\Delta p/p$, where $p$ is the particle momentum and $\Delta p$ is the uncertainty in its transverse component. Since $\Delta p\sim\hbar/\Delta\rho\gg\hbar/\rho$,

\SourcePage{630}{633}
we have $\Delta\theta\gg\hbar/p\rho$, and consequently, in any event,
\begin{equation}
 \theta\gg\frac{\hbar}{\rho mv}.
 \tag{127.8}
\end{equation}
Replacing the angular momentum $m\rho v$ by $\hbar l$, we obtain $\theta l\gg1$, which is the same as the condition $\delta_l\gg1$ (because $\delta_l\sim l\theta$, as follows from (127.3)).

The classical deflection angle can be estimated as the ratio of the transverse momentum increment $\Delta p$ acquired during the collision time $\tau\sim\rho/v$ to the initial momentum $mv$. The force on a particle in the field $U(r)$ at a distance $\rho$ is $F=-dU(\rho)/d\rho$, and hence $\Delta p\sim F\rho/v$, so that $\theta\sim\rho F/mv^2$. Strictly, this estimate is valid only for $\theta\ll1$, but as an order-of-magnitude estimate it may be extended to $\theta\sim1$. Substitution in (127.8) gives the condition for semiclassical scattering:
\begin{equation}
 |F|\rho^2\gg\hbar v.
 \tag{127.9}
\end{equation}
This inequality must hold for all values of $\rho$ for which $|U(\rho)|\lesssim E$ still holds.

If the field $U(r)$ falls off faster than $1/r$, condition (127.9) necessarily fails at sufficiently large $\rho$. Large $\rho$, however, correspond to small $\theta$; hence scattering through sufficiently small angles cannot be classical. If instead the field falls off more slowly than $1/r$, small-angle scattering will be classical. Whether large-angle scattering is also classical in this case depends on how the field behaves at short distances.


For the Coulomb field $U=\alpha/r$, condition (127.9) holds when $\alpha\gg\hbar v$. This is the reverse of the condition under which the Coulomb field may be treated as a perturbation. We shall see, however, that for accidental reasons the quantum theory of scattering in a Coulomb field gives the same result as the classical theory in every case.

\subsection*{Problems}

1. Find the total semiclassical scattering cross-section in a field whose form at sufficiently large distances is $U=\alpha/r^n$, with $n>2$.

\SolutionHeading Bearing in mind that the main contribution comes from phases $\delta_l$ with large $l$, calculate them from (124.1):
\begin{equation}
 \delta_l=-\frac{m\alpha}{\hbar^2}\int_{l/k}^{\infty}\frac{dr}{r^n\sqrt{k^2-l^2/r^2}}
 =-\frac{m\alpha k^{n-2}}{2\hbar^2l^{n-1}}
 \frac{\Gamma(1/2)\Gamma((n-1)/2)}{\Gamma(n/2)}.
 \tag{1}
\end{equation}

\SourcePage{631}{634}
(for evaluation of the integral, cf. Problem 5 of \S~126). Replacing the summation in (123.11) by integration, write
\[
 \sigma=\frac{4\pi}{k^2}\int_0^\infty2l\sin^2\delta_l\,dl.
\]
After the substitution $\delta_l=u$ and an integration by parts with respect to $du$, the integral reduces to a $\Gamma$-function. The result is
\begin{equation}
 \sigma=\frac{2\pi}{n-1}\sin\left[\frac\pi2\left(\frac{n-3}{n-1}\right)\right]
 \Gamma\left(\frac{n-3}{n-1}\right)
 \left[\frac{\sqrt\pi\,\Gamma((n-1)/2)}{\Gamma(n/2)}\right]^{2/(n-1)}
 \left(\frac{m\alpha}{\hbar^2k}\right)^{2/(n-1)}.
 \tag{2}
\end{equation}
(for $n=3$, resolving the indeterminate form gives $\sigma=2\pi^2\alpha m/\hbar v$).

For this result to apply, it is necessary first of all that $l\gg1$ when $\delta_l\sim1$; this gives the inequality
\[
 \frac{m\alpha k^{n-2}}{\hbar^2}\gg1.
\]
Another condition follows from requiring the field $U(r)$ to have the form under consideration already at distances
\[
 r\sim l/k\sim(m\alpha/\hbar^2k)^{1/(n-1)}
\]
(where $l$ is obtained from $\delta_l\sim1$), which make the main contribution to the integral (1). If this form is reached only at distances $r\gg a$, where $a$ is the characteristic size of the field, this gives the condition
\[
 \frac{m\alpha}{\hbar^2ka^{n-1}}\gg1,
\]
which sets an upper limit on the permissible speeds. Recall that in the opposite case, at sufficiently high speeds and subject to $m\alpha/\hbar^2ka^{n-1}\ll1$, one has $\sigma\sim k^{-2}$ (cf. Problem 5 of \S~126).

\ProblemHeading{2.} Determine the distribution of scattering over angle near an extremum of the classical deflection angle $\theta(\rho)$, regarded as a function of the impact parameter $\rho=l/k$.


\SolutionHeading An extremum of $\theta(l)$ at some $l=l_0$ means, by (127.3), that the phase near this point has the form
\[
 \delta_l\approx\delta_{l_0}+\frac{\theta_0}{2}l'+\frac\alpha3l'^3,
\]
where $\theta_0=\theta(l_0)$ and $l'=l-l_0$ (again choosing, for definiteness, the case of the lower sign in (127.3)); the constant $\alpha$ is negative or positive according as $\theta(l)$ has a maximum or a minimum. In place of (127.6), the scattering amplitude is
\[
 |f(\theta)|=\frac1k\left(\frac{l_0}{2\pi\sin\theta_0}\right)^{1/2}
 \left|\int_{-\infty}^{\infty}\exp\left[i\left(\frac{2\alpha}3l'^3-\theta'l'\right)\right]dl'\right|,
\]
Here $\theta'=\theta-\theta_0$. Expressing the integral in terms of the Airy function by (b.3), we finally obtain for the scattering cross section\footnote{Scattering of this type occurs in the theory of the rainbow and is therefore called \emph{rainbow scattering}.}

\[
 d\sigma=\frac{4\pi l_0}{\alpha^{2/3}k^2}
 \Phi^2\left(-\frac{\theta'}{\alpha^{1/3}}\right)d\theta.
\]
